\chapter[Conclusion]{Conclusion}
\label{chap:conclusion}

\section{Summary}
Zeroth-order optimization is prohibitively expensive in high dimensions, typically needing $d$ times more function evaluations than gradient descent. When perturbations are Gaussian, parallelizing function evaluations across $r$ workers typically reduces the number of iterations by a factor of $r$. EGGROLL replaces these Gaussian perturbations with structured low-rank matrices $ab^\top$ for Gaussian vectors $a,b$~\citep{sarkar2025evolution}. Suddenly, zeroth-order algorithms were practical. Yet it remained open whether these algorithms provably converge.

\paragraph{EGGROLL Convergence.}
In Chapter~\ref{chap:fosp-eggroll}, we answered the above question in the affirmative, proving that, after a judicious choice of parameters, EGGROLL converges to a first-order stationary point for any $L$-smooth (possibly nonconvex) function $f$. The principal theoretical contribution was Theorem~\ref{thm:error_concentration}, showing the typical norm of the gradient estimator using low-rank perturbations is the same as with Gaussian perturbations up to irrelevant log factors, vastly generalizing Proposition 2 of~\citet{ren2023escaping}. Having established the theorem, we can prove that EGGROLL converges.

\paragraph{Zeroth-Order Mirror Descent.}
Perturbed Gradient Descent was shown to rapidly escape saddle points with only a logarithmic penalty on dimension~\citep{jin2017escape}. \citet{ren2023escaping} showed that perturbed Zeroth-Order Gradient Descent escapes saddle points in time $O(1/\ve^{5/2})$. Rather than just generalize zeroth-order perturbed gradient descent to the low-rank case, we seek a better zeroth-order rate such as the $O(1/\ve^{7/4})$ demonstrated in the first-order methods of~\citet{carmon2018accelerated} and~\citet{jin2017accelerated}. \citet{carmon2018accelerated} design a generalizable method which relies on two subroutines: accelerated gradient descent and finding an eigenvector corresponding to the smallest Hessian eigenvalue. In pursuit of a zeroth-order accelerated gradient descent, we first showed that zeroth-order mirror descent converges in Chapter~\ref{chap:zo-mirror-descent} with the typical $d$ zeroth-order penalty for quadratic convex functions.

\paragraph{Zeroth-Order Accelerated Gradient Descent.} Having established the theoretical tools needed for proving convergence of zeroth-order mirror descent, we use the linear mixing framework introduced by~\citet{allenZhu2016linear} to design a zeroth-order accelerated gradient descent algorithm which converges with the typical $d$ zeroth-order penalty, with high probability, in a single run. We precisely quantify how an inexact gradient oracle propagates error through the algorithm by bounding the martingale term by other $\Theta(1)$ terms already present. We believe the techniques developed in this chapter are generalizable to other problems that need to bound similar terms.

\paragraph{Impact.} Armed with the ability to bound mean-zero martingale error terms and the typical scale of the EGGROLL gradient estimator with high probability, many high-probability proofs in zeroth-order descent can be generalized to this low-rank perturbation setting. Low-rank directions are especially suitable for matrix-valued neural network parameters, substantially reducing memory requirements from $O(rdm)$ to $O(r(d+m))$, allowing for the use of zeroth-order algorithms in practice~\citep{sarkar2025evolution}. In Chapters~\ref{chap:fosp-eggroll} and~\ref{chap:linear-mixing} we demonstrate that parallelizing function evaluations across $r$ workers yields an $r$-fold reduction in the number of iterations. As investment in compute continues to grow, these results suggest zeroth-order methods could become a useful complement to gradient-based reinforcement-learning methods such as PPO and GRPO.

\section{Limitations}
\paragraph{Optimizing $\sigma$.}
Here, we have not attempted to optimize the smoothing radius $\sigma$. In the deterministic case, larger $\sigma$ allows for more numerical stability. In the stochastic case with additive noise $\mu$, one has to balance the bias from the smoothing radius $\sigma$ with an error term of order $\mu/\sigma$ introduced by the gradient estimator. Therefore, increasing $\sigma$ reduces noise amplification but increases smoothing bias. For example, \citet{he2024riemannian} highlight that their method permits a larger smoothing radius.

\paragraph{Logarithms.}
We have ignored polylogarithmic factors throughout this work. It might be possible to remove logarithmic dependence on $\ve$ in Chapter~\ref{chap:linear-mixing}, but we have not attempted it here. Although logarithmic factors are commonly suppressed in analyses of randomized algorithms, some work has focused on removing this~\citep{li2022last}.

\paragraph{Function-Evaluation Complexity.}
While we have shown that parallelizing function evaluations across $r$ workers reduces wall-clock time by a factor of $r$, the total number of function evaluations is still $d$ times larger than gradient descent. To the best of our knowledge, no existing method achieves a sublinear dependence on $d$ for the number of function evaluations for optimizing convex functions in the zeroth-order setting.

\paragraph{Logarithmic Lower Bound on $r$.} In Chapter~\ref{chap:linear-mixing}, we assume a lower bound $r=\widetilde{\Omega}(1)$. Significant work in~\citet{ren2023escaping} and Chapter~\ref{chap:fosp-eggroll} went into removing this assumption. However, we expect it to be necessary for this zeroth-order linear coupling framework since we reset $z_0$ back to its original value every restart. \textit{A priori}, $z_0$ and $x_0$ could drift very far from each other, preventing us from establishing that gradients change by at most a factor of 4 over an interval of log-length unlike Lemma~\ref{lem:half-two-gradient}.

\paragraph{$L$-Smoothness Assumption.} All our results depend on an $L$-smoothness assumption. While for $f \in \mathcal C^1(\R^d)$ this is more general than a Lipschitz assumption, which frequently does not hold for ML models, it is still restrictive. In the EGGROLL paper~\citep{sarkar2025evolution}, they assume only a local $\alpha$-H\"older condition on gradients. 

\section{Future Work}
\paragraph{Experiments.}
This work is purely theoretical, so no experiments were performed. In future work, significant experimental investigation should tune $\alpha$ and $\eta$ for zeroth-order linear mixing (Chapter~\ref{chap:linear-mixing}) along with $\eta$ and $\sigma$ for zeroth-order convergence to a FOSP (Chapter~\ref{chap:fosp-eggroll}). For large ML models, we typically only use batched losses, which makes function evaluations stochastic, not deterministic as we have assumed here. This might necessitate a variance-reduction technique, or adaptive choices for $\eta$ and $\alpha$ such as the conditions of~\citet{robbins1951stochastic}. 

\paragraph{Anisotropic Perturbations.}
We have assumed throughout that perturbations are sampled with covariance $I$. However, there has been some work on using learned anisotropic perturbations \citep{seung2026lowrank}. The best anisotropic estimator would have covariance $\grad f(x) \grad f(x)^\top$ at $x$, since we could immediately recover the gradient direction. Perhaps remembering a time-average of past gradients ($\tau < t$) could inform the gradient at time $t$, thereby reducing the variance of the estimator. There has been some preliminary work in this direction~\citep{liu2020selfguided,meier2019improving}.

\paragraph{Generalizing Existing Zeroth-Order Methods.} Armed with Theorem~\ref{thm:error_concentration} and its consequences, we suspect many existing ZO methods can be generalized to the case of low-rank perturbations. We have pursued this for~\citet{ren2023escaping}, but we suspect there is more to do here.

\section{Concluding Remarks}
Low-rank perturbations are a powerful tool for making zeroth-order algorithms practical for neural networks. This work shows this low-rank structure still permits rigorous high-probability guarantees, rather than being merely an engineering device. We showed how zeroth-order accelerated convex and nonconvex optimization can be analyzed with high probability, quantifying propagated error through a mean-zero martingale error term. We showed how the population size $r$ trades parallel compute for a decrease in wall-clock time. While other fine-tuning methods like GRPO and PPO are only partially understood theoretically, we establish that EGGROLL converges to a FOSP in the deterministic $L$-smooth setting. Low-rank perturbations may therefore make zeroth-order optimization practical at scale. This dissertation establishes a theoretical foundation for their use and opens several directions for future work.